\subsection{Linear Equations in Three Variables} 
What does an equation in three variables ``look like''?
\begin{Example}
Sketch the graph of $x+y+z=2$
\begin{itemize}
\item Remember, the graph of an equation is a picture of all the \makeblank[1in] to the equation.
\item Since there are \makeblanksmall variables in this equation, its solution is an ordered \emph{triple}, that is, \makeblanksmall values $(x,y,z)$.
\item To graph this, we need \makeblanksmall axes. We draw the $z$-axis on a \makeblank[1in], but imagine that it ``comes out of the screen'' toward you.
\item It is easy to see that \blanktrip, \blanktrip, and \blanktrip\ are solutions to this equation.
\end{itemize}
\begin{icpic}
\useasboundingbox (-5.5,-5.5) rectangle (5.5,5.5);
\draw (-5,5) node[anchor=north east] {$x+y+z=2$};
\setgraphsquare{5}
\AxesGen{\small}{\scriptsize}
\zAxis{\small}{\scriptsize}
\end{icpic}
\end{Example}
The graph of a linear equation in three variables is a \makeblank[1in] in space. When a system of three equations has a unique solution, we have three \makeblank[1in] intersecting at a \makeblank[1in]:
\begin{icpic}[x=2\linespace,y=2\linespace]
\draw[pca!50!pcb,plane,dashed] (\project{-1}{0}{0}) -- (0,0);
\draw[pca!50!pcc,plane,dashed] (\project{0}{-1}{0}) -- (0,0);
\draw[pcb!50!pcc,plane,dashed] (\project{0}{0}{-1}) -- (0,0);
\filldraw[pca,plane]  (\project{-1}{-1}{0}) -- 
                      (\project{-1}{1}{0}) --
										  (\project{1}{1}{0}) --
											(\project{1}{-1}{0}) --
											cycle;
\filldraw[pcb,plane]  (\project{-1}{0}{-1}) --
                      (\project{-1}{0}{1}) --
											(\project{1}{0}{1}) --
											(\project{1}{0}{-1}) -- 
											cycle;
\filldraw[pcc,plane]  (\project{0}{-1}{-1}) -- 
                      (\project{0}{-1}{1}) -- 
											(\project{0}{1}{1}) -- 
											(\project{0}{1}{-1}) -- 
											cycle;
\draw[pca!50!pcb,plane] (\project{1}{0}{0}) -- (0,0);
\draw[pca!50!pcc,plane] (\project{0}{1}{0}) -- (0,0);
\draw[pcb!50!pcc,plane] (\project{0}{0}{1}) -- (0,0);
\end{icpic}
What could go wrong?